\subsection{COMPARISON OF TOPOLOGIES}

Theorem 2 of no. 1 shows that we may take the continuous mappings as morphisms of topological structures (Set Theory, Chapter IV, § 2, no. 1); from now on, we shall assume that we have made this choice of morphisms. In accordance with the general definitions (Set Theory, Chapter IV, § 2, no. 2), this allows us to define an ordering on the set of topologies on a given set X:

DEFINITION 3. Given two topologies \(\mathfrak{C}_1\) , \(\mathfrak{C}_2\) on the same set X, we say that \(\mathfrak{C}_1\) is finer than \(\mathfrak{C}_2\) (and that \(\mathfrak{C}_2\) is coarser than \(\mathfrak{C}_1\) ) if, denoting by \(X_i\) the set X with the topology \(\mathfrak{C}_i\) ( \(i = 1, 2\) ), the identity mapping \(X_1 \to X_2\) is continuous. If in addition \(\mathfrak{C}_1 \neq \mathfrak{C}_2\) , we say that \(\mathfrak{C}_1\) is strictly finer than \(\mathfrak{C}_2\) (and that \(\mathfrak{C}_2\) is strictly coarser than \(\mathfrak{C}_1\) ).

Two topologies, one of which is finer than the other, are said to be comparable.

The criteria for a mapping to be continuous (no. 1, Definition 1 and Theorem 1) give the following proposition:

PROPOSITION 3. Given two topologies \(\mathfrak{G}_1\) , \(\mathfrak{G}_2\) on a set \(X\) , the following statements are equivalent:

\begin{enumerate}
\def\labelenumi{\alph{enumi})}
\item
  \(\mathfrak{C}_1\) is finer than \(\mathfrak{C}_2\) .
\item
  For each \(x \in \mathbf{X}\) , each neighbourhood of \(x\) for \(\mathfrak{C}_2\) is a neighbourhood of \(x\) for \(\mathfrak{C}_1\) .
\item
  For each subset A of X, the closure of A in the topology \(C_{1}\) is contained in the closure of A in the topology \(C_{2}\) .
\item
  Every subset of X which is closed in \(\mathfrak{C}_{2}\) is closed in \(\mathfrak{C}_{1}\) .
\item
  Every subset of X which is open in \(\mathfrak{C}_{2}\) is open in \(\mathfrak{C}_{1}\) .
\end{enumerate}

Example. * In Hilbert space H consisting of sequences \(x = (x_n)\) of real numbers such that

\[
\| x \| ^ {2} = \sum_ {n = 0} ^ {\infty} x _ {n} ^ {2} <   + \infty ,
\]

the neighbourhoods of a point \(x_0\) in the strong topology on H are the sets which contain an open ball \(\|x - x_0\| < a\) centred at \(x_0\) ; the neighbourhoods of \(x_0\) in the weak topology on H are the sets containing a set defined by relation of the form \(\sup_{1 \leqslant i \leqslant n} |(x - x_0|a_i)| \leqslant 1\) , where the \(a_i\) are points of H and

\[
(x | y) = \sum_ {n = 0} ^ {\infty} x _ {n} y _ {n}
\]

if \(x = (x_n)\) and \(y = (y_n)\) . Now if \(\beta = \sup_{1 \leqslant i \leqslant n} \|a_i\|\) , the relation \(\|x - x_0\| \leqslant 1/\beta\) implies that \(|(x - x_0|a_i)| \leqslant \|x - x_0\| \cdot \|a_i\| \leqslant 1\) for \(1 \leqslant i \leqslant n\) ; hence the strong topology on H is finer than the weak topology. On the other hand, given any finite family \((a_i)_{1 \leqslant i \leqslant n}\) of points of H, there are points \(x\) in H such that \((x - x_0|a_i) = 0\) for \(1 \leqslant i \leqslant n\) and such that \(\|x - x_0\|\) is arbitrarily large; this shows that the strong topology is strictly finer than the weak topology. *

Remarks. 1) In the ordered set of all topologies on a set X, the topology in which the only open sets are \(\varnothing\) and X is the coarsest and the discrete topology is the finest.

\begin{enumerate}
\def\labelenumi{\arabic{enumi})}
\setcounter{enumi}{1}
\item
  The finer the topology, the more open sets, closed sets and neighbourhoods; the finer the topology, the smaller (resp. the larger) the closure (resp. the interior) of a set; the finer the topology, the fewer dense sets.
\item
  If \(f: \mathbf{X} \to \mathbf{X}'\) is a continuous mapping, it remains continuous if the topology of \(\mathbf{X}\) is replaced by a finer topology and the topology of \(\mathbf{X}'\) is replaced by a coarser topology (no. 1, Theorem 2). In other words, the finer the topology of \(\mathbf{X}\) and the coarser the topology of \(\mathbf{X}'\) , the more continuous mappings there are of \(\mathbf{X}\) into \(\mathbf{X}'\) .
\end{enumerate}

